\chapter{Literature Review}
\label{ch:literature-review}
This chapter asks which approaches already recover a sharp fluorescence image through optical aberrations, and what each requires. We compare them along five requirements: a known or estimated PSF, additional optical hardware, additional captures, coherent illumination, and whether one correction holds across the field. The sections follow where the correction takes place: after capture, by deconvolution with a PSF (Section 2.1); at the optics, with adaptive optics (2.2); through additional acquisitions (2.3); and from light-field data (2.4).

\section{PSF-based computational correction}
\label{sec:lit-psf-correction}
The PSF is the image that an optical system forms of a point source. Most deconvolution algorithms model the microscope as linear and shift-invariant, so that the image is the scene convolved with a single PSF, and they estimate the scene by inverting that convolution~\cite{sarder2006deconvolutionMethods,sibarita2005deconvolutionMicroscopy}. Real samples break this model. A refractive-index mismatch between the immersion and mounting media introduces spherical aberration that grows with imaging depth~\cite{preza2004depthVariantRestoration}, and refractive-index variation inside the specimen distorts different regions differently~\cite{kam2001computationalAdaptiveOptics}, so in practice the PSF is not shift-invariant~\cite{sarder2006deconvolutionMethods}. The PSF can be computed from a model, measured, or estimated blindly, and each route has its own limitation. A theoretical PSF misses the irregularities of the particular microscope, while a measured PSF carries noise and describes only the conditions under which it was recorded~\cite{hanser2004phaseRetrievedPupil}. Blind estimation recovers both the scene and the PSF from the same data, an ill-posed and underdetermined problem; it is needed because the sample itself alters the PSF~\cite{kenig2010blindDeconvolutionLearning}. Whichever route is taken, an inaccurate PSF produces artifacts in the restored image~\cite{sibarita2005deconvolutionMicroscopy,mcnally1999deconvolutionMicroscopy3d,preza2004depthVariantRestoration}. The inversion has costs of its own. Direct inversion amplifies noise~\cite{sarder2006deconvolutionMethods}. Statistical methods model the noise and are more robust~\cite{sibarita2005deconvolutionMicroscopy,mcnally1999deconvolutionMicroscopy3d}, yet Richardson--Lucy, the maximum-likelihood method for Poisson noise, still amplifies it unless stopped early or regularized~\cite{dey2006richardsonLucyTv}. Iterative and blind methods are computationally heavy; one blind scheme took about two hours for a $128^3$ volume~\cite{kenig2010blindDeconvolutionLearning}. A spatially varying aberration needs a separate PSF for each region. Kam et al.\ predicted such PSFs by ray tracing through a refractive-index map of the specimen and deconvolved region by region, but found full point-by-point space-variant deconvolution computationally intractable for 3D data, and demonstrated the approach on fluorescent beads only~\cite{kam2001computationalAdaptiveOptics}. PSF-based correction therefore requires a known or estimated PSF for every region of the field, a requirement that the method in this thesis avoids.

\section{Light-field microscopy}
\label{sec:lit-light-field-microscopy}
Light-field microscopy (LFM) records spatial information together with angular information in a single exposure, and this snapshot capability has made it important in high-speed biological imaging, including volumetric demonstrations such as whole-animal neuronal activity imaging \cite{prevedel2014wholeAnimalLfm}. At the same time, LFM performance depends on the optical geometry, calibration, and reconstruction model, and reconstruction artifacts can arise when the captured information is undersampled or poorly modeled \cite{levoy2009recordingControlling4d,broxton2013waveOpticsLfm,stefanoiu2019artifactFreeDeconvolution}. In image-plane microlens-array LFM, the reconstruction model can also involve a position-dependent point-spread function (PSF) that maps a three-dimensional scene to the two-dimensional sensor measurement, increasing computational cost for volumetric reconstruction \cite{broxton2013waveOpticsLfm,guo2019fourierLightField}. These limitations make the optical architecture used to encode the light field central to both the information that can be recovered and the cost of interpreting it. Fourier light-field microscopy (FLFM), which places the microlens array in a pupil-conjugate plane, addresses these limitations. In the implementation of Guo et al. \cite{guo2019fourierLightField}, this arrangement reduced sampling-related artifacts and reconstruction time relative to conventional LFM; these benefits depend on the optical design, calibration, and reconstruction model rather than following unconditionally.
